
\documentclass{article}
\usepackage{colortbl}
\usepackage{makecell}
\usepackage{multirow}
\usepackage{supertabular}

\begin{document}

\newcounter{utterance}

\centering \large Interaction Transcript for game `bbh', experiment `formal\_fallacies', episode 99 with Qwen3.6{-}27B.
\vspace{24pt}

{ \footnotesize  \setcounter{utterance}{1}
\setlength{\tabcolsep}{0pt}
\begin{supertabular}{c@{$\;$}|p{.15\linewidth}@{}p{.15\linewidth}p{.15\linewidth}p{.15\linewidth}p{.15\linewidth}p{.15\linewidth}}
    \# & \multicolumn{2}{c}{Player} && \multicolumn{2}{c}{Game Master} \\
    \hline

    \theutterance \stepcounter{utterance}  
    & & & \multicolumn{4}{p{0.6\linewidth}}{
        \cellcolor[rgb]{0.9,0.9,0.9}{
            \makecell[{{p{\linewidth}}}]{
                \texttt{\tiny{[P1$\langle$GM]}}
                \texttt{Distinguish deductively valid arguments from formal fallacies.} \\
\\ 
\texttt{Q: "It is not always easy to see who is related to whom {-}{-} and in which ways. The following argument pertains to this question: To begin with, Lesley is a close friend of Fernando. Moreover, being a close friend of Fernando or a schoolmate of Lowell is sufficient for being a great{-}grandfather of Leroy. It follows that Lesley is a great{-}grandfather of Leroy."} \\
\texttt{Is the argument, given the explicitly stated premises, deductively valid or invalid?} \\
\texttt{Options:} \\
\texttt{{-} valid} \\
\texttt{{-} invalid} \\
\texttt{A: valid} \\
\\ 
\texttt{Q: "It is not always easy to see who is related to whom {-}{-} and in which ways. The following argument pertains to this question: Whoever is not a great{-}grandfather of Clyde is a stepbrother of Brian. Being an ancestor of Dana is sufficient for not being a great{-}grandfather of Clyde. We may conclude: Everyone who is an ancestor of Dana is a stepbrother of Brian, too."} \\
\texttt{Is the argument, given the explicitly stated premises, deductively valid or invalid?} \\
\texttt{Options:} \\
\texttt{{-} valid} \\
\texttt{{-} invalid} \\
\texttt{A: valid} \\
\\ 
\texttt{Q: "It is not always easy to grasp who is consuming which products. The following argument pertains to this question: Every infrequent user of Paul Mitchell shampoo is either a rare consumer of Nioxin shampoo or a loyal buyer of Caress soap, or both. No regular consumer of Lush soap is a rare consumer of Nioxin shampoo and, in the same time, a loyal buyer of Caress soap. It follows that whoever is an infrequent user of Paul Mitchell shampoo is not a regular consumer of Lush soap."} \\
\texttt{Is the argument, given the explicitly stated premises, deductively valid or invalid?} \\
\texttt{Options:} \\
\texttt{{-} valid} \\
\texttt{{-} invalid} \\
\texttt{A: invalid} \\
\\ 
\texttt{Q: "Here comes a perfectly valid argument: To start with, everyone who is both a rare consumer of Pantene shampoo and an occasional purchaser of Herbal Essences shampoo is a frequent consumer of Tom's of Maine soap or a regular user of Mysore soap. Now, every frequent consumer of Tom's of Maine soap is neither a rare consumer of Pantene shampoo nor an occasional purchaser of Herbal Essences shampoo. We may conclude that everyone who is not both a rare consumer of Pantene shampoo and an occasional purchaser of Herbal Essences shampoo is however a regular user of Mysore soap."} \\
\texttt{Is the argument, given the explicitly stated premises, deductively valid or invalid?} \\
\texttt{Options:} \\
\texttt{{-} valid } \\
\texttt{{-} invalid} \\
\texttt{A:} \\
            }
        }
    }
     \\ \\

    \theutterance \stepcounter{utterance}  
    & \multicolumn{4}{p{0.6\linewidth}}{
        \cellcolor[rgb]{0.561,0.737,0.561}{
            \makecell[{{p{\linewidth}}}]{
                \texttt{\tiny{[P1$\rangle$GM]}}
                \texttt{To determine whether the argument is deductively valid, let's break it down into logical components.} \\
\\ 
\texttt{**1. Symbolize the Premises and Conclusion:**} \\
\\ 
\texttt{Let:} \\
\texttt{*   \$P\$: Rare consumer of Pantene shampoo} \\
\texttt{*   \$H\$: Occasional purchaser of Herbal Essences shampoo} \\
\texttt{*   \$T\$: Frequent consumer of Tom's of Maine soap} \\
\texttt{*   \$M\$: Regular user of Mysore soap} \\
\\ 
\texttt{**Premise 1:** "Everyone who is both a rare consumer of Pantene shampoo and an occasional purchaser of Herbal Essences shampoo is a frequent consumer of Tom's of Maine soap or a regular user of Mysore soap."} \\
\texttt{*   Logic: \$(P \textbackslash{}land H) \textbackslash{}rightarrow (T \textbackslash{}lor M)\$} \\
\\ 
\texttt{**Premise 2:** "Every frequent consumer of Tom's of Maine soap is neither a rare consumer of Pantene shampoo nor an occasional purchaser of Herbal Essences shampoo."} \\
\texttt{*   Logic: \$T \textbackslash{}rightarrow \textbackslash{}neg(P \textbackslash{}land H)\$} \\
\texttt{    *   Note: \$\textbackslash{}neg(P \textbackslash{}land H)\$ is equivalent to \$\textbackslash{}neg P \textbackslash{}lor \textbackslash{}neg H\$.} \\
\\ 
\texttt{**Conclusion:** "Everyone who is not both a rare consumer of Pantene shampoo and an occasional purchaser of Herbal Essences shampoo is however a regular user of Mysore soap."} \\
\texttt{*   Logic: \$\textbackslash{}neg(P \textbackslash{}land H) \textbackslash{}rightarrow M\$} \\
\\ 
\texttt{**2. Analyze the Validity:**} \\
\\ 
\texttt{We need to check if the conclusion} \\
            }
        }
    }
    & & \\ \\

    \theutterance \stepcounter{utterance}  
    & & & \multicolumn{2}{p{0.3\linewidth}}{
        \cellcolor[rgb]{0.9,0.9,0.9}{
            \makecell[{{p{\linewidth}}}]{
                \texttt{\tiny{[GM$|$GM]}}
                \texttt{Target: invalid} \\
            }
        }
    }
    & & \\ \\

    \theutterance \stepcounter{utterance}  
    & & & \multicolumn{2}{p{0.3\linewidth}}{
        \cellcolor[rgb]{0.9,0.9,0.9}{
            \makecell[{{p{\linewidth}}}]{
                \texttt{\tiny{[GM$|$GM]}}
                \texttt{game\_result = LOSE} \\
            }
        }
    }
    & & \\ \\

\end{supertabular}
}

\end{document}
